\section{Convexity and affine manifolds}

\subsection{Classical convexity}
When~$M=\bbR^n$, there is a notion of \emph{strict convexity}
of a function~$u$, which is the condition
that the Hessian quadratic form~$H(u)$ be positive definite, where
\begin{equation}\label{eq: flat_Hessian}
H(u) = {\frac{\partial^2 u}{\partial x^i\partial x^j}}
\,\d x^i\otimes\d x^j
\end{equation}
and where $x^1,\ldots,x^n$
are the usual affine linear coordinates in~$\bbR^n$.
Note that strict convexity is an affine-invariant notion on~$\bbR^n$.

Motivated by applications in economics,
Ekeland and Nirenberg~\cite{MR2023129}
asked what further conditions one must impose on an~$\omega\in\Omega^1(\bbR^n)$
satisfying~\eqref{eq: omegaPfaff_k} and~\eqref{eq: omega_nondegenerate}
in order to know that one can choose the functions~$u^j$
and $a_j$ in the representation~\eqref{eq: PfaffDarboux_normalform}
so that the $u^j$ be strictly convex and the $a_j$ be positive.
It is not hard to show, by example, that \emph{some}
further condition on~$\omega$ is necessary
to guarantee the existence of such a convex representation.
(See the discussion at the beginning of Section~\ref{ssec: poscond}.)

They showed that two earlier articles~\cite{MR1655519,MR1652709}
claiming to provide such necessary and sufficient conditions were flawed
(indeed, they exhibited counterexamples to the claims of these articles)
and then produced their own condition,
which they showed to be necessary and sufficient.

In this note, I will show that their main result,
properly formulated, holds good on an $n$-manifold~$M$
endowed with a torsion-free affine connection,
not just on~$\bbR^n$ endowed with the (flat) affine connection
it inherits as a vector space.

\subsection{Affine connections and convexity}

Let~$\nabla$ be a torsion-free affine connection on an $n$-manifold~$M^n$,
i.e., $\nabla$ is a first-order, linear differential operator
\begin{equation*}
\nabla \colon \ \Omega^1(M)\to \Omega^1(M)\otimes \Omega^1(M)
\end{equation*}
that obeys the Leibnitz rule
\begin{equation*}
\nabla (f\eta) = \d f\otimes\eta + f \nabla(\eta)
\end{equation*}
for all smooth functions~$f$ on~$M$ and smooth $1$-forms~$\eta$ on~$M$.
The assumption that~$\nabla$ be torsion-free
is the condition that the associated (second-order)
\emph{Hessian operator}~$H(u)=\nabla(\d u)$
be a symmetric $(0,2)$-tensor for each smooth function~$u$ on~$M$.

A smooth function~$u$ on~$M$ is said to be {\it strictly $\nabla$-convex}
if, as a quadratic form,~$H(u)$ is positive definite at every point of~$M$.

When~$M=\bbR^n$ and~$\nabla$ is the standard (flat) connection,
satisfying~$\nabla(\d x^i) = 0$ for all of the coordinate functions~$x^i$,
then~$H(u)$ is the usual Hessian tensor~\eqref{eq: flat_Hessian},
and this notion of convexity is simply the classical one.

In the more general case, when $x = (x^i)\colon U\to \bbR^n$ is a local
coordinate chart, one has
\begin{equation*}
H\big(x^k\big) = \nabla\big(\d x^k\big) = \Gamma^k_{ij}\,\d x^i\otimes\d x^j,
\end{equation*}
where $\Gamma^k_{ij}=\Gamma^k_{ji}\in C^\infty(U)$
are the coefficients of the connection $\nabla$ relative
to the coordinate chart~$x=(x^i)$.
The general coordinate formula for~$H$ then becomes
\begin{equation*}
H(u) = \left(\frac{\partial^2 u}{\partial x^i\partial x^j}
+ \Gamma^k_{ij} \frac{\partial u}{\partial x^k}\right)
\,\d x^i\otimes\d x^j.
\end{equation*}
Thus, $\nabla$-convexity of $u$ is expressible in terms of a condition
on the $2$-jet of~$u$, slightly more general than the condition for
classical convexity.

Adopting the usual conventions
\begin{equation*}%\label{eq: altsym_decomp}
\alpha\w\beta
= \tfrac12(\alpha\otimes\beta-\beta\otimes\alpha), \qquad
\alpha\circ\beta
= \tfrac12(\alpha\otimes\beta+\beta\otimes\alpha),
\end{equation*}
one sees that, for a $1$-form~$\omega$ of the form
\begin{equation}\label{eq: P-Drepresentation}
\omega = a_1\,\d u^1 + \cdots + a_k\,\d u^k,
\end{equation}
one has (using the summation convention)
\begin{align*}
\nabla\omega
&= \d a_i\otimes\d u^i + a_i \nabla(u^i) = \d a_i\otimes\d u^i + a_i H(u^i)\\
&= \d a_i\w\d u^i + \d a_i\circ\d u^i + a_i H(u^i)\\
&= \d\omega + \Ss^\nabla\omega,
\end{align*}
where I have introduced
the notation~$\Ss^\nabla\omega$ to denote the {\it symmetrization\/}
of~$\nabla(\omega)$. Thus,~$\Ss^\nabla\omega
= \nabla\omega- \d\omega$
is a well-defined quadratic form on~$M$.
(Of course, the linear, first-order differential operator~$\Ss^\nabla$
depends on~$\nabla$.)


\subsection{A positivity condition}\label{ssec: poscond}
The quadratic form $S^\nabla\omega$ provides some insight
into the question of whether a $\nabla$-convex Pfaff--Darboux
representation of $\omega$ is possible.

\begin{Proposition}\label{prop: ness cond}
Suppose that $\omega\in\Omega^1(M)$ satisfies~\eqref{eq: omegaPfaff_k}
and~\eqref{eq: omega_nondegenerate}. If there exist positive
functions $a_i$ and $\nabla$-convex functions $u^i$ for $1\le i\le k$
such that~\eqref{eq: P-Drepresentation} holds, then $S^\nabla\omega$
is positive definite on the leaves of the foliation $\cL$ defined by
$\d u^1 = \d u^2 = \cdots = \d u^k = 0$.
\end{Proposition}

\begin{proof}
Since $\Ss^\nabla\omega =\d a_i\circ \d u^i + a_i H(u^i)$,
it follows that, when restricted to the plane field ${N\subset TM}$
defined by~$\d u^1 = \d u^2 = \cdots = \d u^k = 0$,
the terms~$\d a_i\circ \d u^i$ in~$\Ss^\nabla\omega$
vanish. Thus, $S^\nabla\omega = a_i H(u^i)$
as quadratic forms on $N$. By the positivity of the $a_i$
and the $\nabla$-convexity of the $u^i$, it follows that
$S^\nabla\omega$ is positive definite on~$N$.\footnote{It is worth pointing out that the same conclusion
about the positive definiteness of $\Ss^\nabla\omega$ on the leaves of~$\cL$
would have followed if one had merely assumed that each $u^i$
be only `strictly $\nabla$-\emph{quasi}-convex',
i.e., that $\d u^i$ be nonvanishing and $H(u^i)$ be positive definite
when restricted to the hyperplane field~$\d u^i = 0$.
Compare~\cite[Lemma~1]{MR2023129},
and the preceding discussion about their Problem~2.}
\end{proof}

This proposition provides necessary condition for the existence
of a convex Pfaff--Darboux representation.

\begin{Example}[an obstructed example]\label{ex: Obstruction}
Let $M=\mathbb{R}^n$ with standard coordinates $x=(x^i)$,
and let $\nabla$ be the (flat, torsion-free) connection
such that $\nabla(\d x^i)=0$ for $1\le i\le n$.
Let $c_i$, $f_{ij}=-f_{ji}$, and $g_{ij}=g_{ji}$
be constants and consider the $1$-form
\[
\omega = \bigl(c_i + (f_{ij}+g_{ij})x^j)\,\d x^i.
\]
Then $\d\omega= f_{ij}\,\d x^j\w\d x^i = - f_{ij}\,\d x^i\w\d x^j$
and $\Ss^\nabla\omega = g_{ij}\,\d x^i\circ\d x^j$.

Now assume that the skew-symmetric matrix $f = (f_{ij})$
has rank $2(k{-}1)<n$, so that $(\d\omega)^{k-1}\not=0$
but $(\d\omega)^{k}=0$.
Then, for generic choice of the constants $c_i$,
$\omega\w(\d\omega)^{k-1}$ will be nonvanishing on an open neighborhood $U\subset\mathbb{R}^n$ of the origin $0\in\mathbb{R}^n$,
in which case, $\omega$ will satisfy the hypotheses~\eqref{eq: omegaPfaff_k}
and~\eqref{eq: omega_nondegenerate} on~$U$.

If the symmetric matrix $g = (g_{ij})$
does not have at least $n{-}k$ positive eigenvalues,
then~$\Ss^\nabla\omega$ cannot be positive definite on
any $(n{-}k)$-dimensional subbundle $N\subset TU$, and, hence,
by Proposition~\ref{prop: ness cond}, $\omega$ cannot have a convex Pfaff--Darboux
representation in any open subset of $U$.
\end{Example}

It turns out that this necessary condition
for a local `convex' Pfaff--Darboux representation compatible
with an $\omega$-Legendrian foliation~$\cL$ is also sufficient.

\begin{Theorem}\label{thm: main}
Suppose $\nabla$ be a torsion-free affine
connection on~$M$, that $\omega\in\Omega^1(M)$
satisfy~\eqref{eq: omegaPfaff_k} and~\eqref{eq: omega_nondegenerate}
for some~$k>0$, and that $\cL$ be an $\omega$-Legendrian
foliation on~$M$ with the property that $\Ss^\nabla\omega$
pulls back to each leaf of~$\cL$ to be positive definite.
Then each $m\in M$ has an open neighborhood~$U\subset M$
on which there exist strictly $\nabla$-convex functions~$u^1,\ldots,u^k$
that are constant on the leaves of~$\cL$ in~$U$
and positive functions~$a_1,\ldots, a_k$ such that
\begin{equation*}
U^*\omega = a_1\,\d u^1 + \cdots + a_k\,\d u^k.
\end{equation*}
\end{Theorem}

